% Slide "Several cores: speedup over plain scalar", right panel: f64. Speedup over the plain scalar
% build of the same compiler at the same core count, lid-driven cavity (SweepgenLdcCPU), one 64^3
% block per MPI process, best of 3 runs of 200 steps. One session (board, scripts/presentation4/all.sh,
% logs docs/presentation4/board-logs/board/ldc-{gcc,clang}-{scalar,autovec,rvv}-f64-{1,4,8}core.log),
% printed by
%   python3 scripts/presentation4/slide_numbers.py --tex docs/presentation4/board-logs/board
% Bars and numbers: GCC 14.2. Diamonds: Clang 18.1, against its own plain scalar (plain scalar = 1
% for both compilers, so no diamond there). Series as figures/mlups-multicore.tex.
% Styles mc, pl, bs, av, rv, sp, cl, cllegend: preamble of slides.tex. Same ymax as the f32 panel.
% The legend is stored as mcspeeduplegend and placed under both panels by \ref in slides.tex.
% presentation5: labels with two digits (notes5 F8).
\begin{tikzpicture}
  \begin{axis}[mc, ymax=7.5, title={f64},
      legend to name=mcspeeduplegend, legend columns=-1,
      legend style={font=\tiny,draw=psgray!50,/tikz/every even column/.append style={column sep=5pt}}]
    \addplot[pl] coordinates {(1,1)[1.0] (4,1)[1.0] (8,1)[1.0]};
    \addplot[bs] coordinates {(1,1.585)[1.6] (4,1.522)[1.5] (8,1.469)[1.5]};
    \addplot[av] coordinates {(1,1.566)[1.6] (4,1.524)[1.5] (8,1.476)[1.5]};
    \addplot[rv] coordinates {(1,4.230)[4.2] (4,2.917)[2.9] (8,1.993)[2.0]};
    \addplot[sp] coordinates {(1,5.123)[5.1] (4,2.968)[3.0] (8,2.004)[2.0]};
    \addplot[cl,xshift=-6pt] coordinates {(1,1.773) (4,1.619) (8,1.510)};
    \addplot[cl,xshift=0pt] coordinates {(1,1.768) (4,1.648) (8,1.524)};
    \addplot[cl,xshift=6pt] coordinates {(1,4.061) (4,2.779) (8,1.931)};
    \addplot[cl,xshift=12pt,forget plot=false,cllegend] coordinates {(1,4.959) (4,2.887) (8,1.976)};
    \legend{plain scalar, best scalar, auto-vec.\ in place, RVV in place + prefetch, RVV split loop,
      Clang}
  \end{axis}
\end{tikzpicture}
